\section{The first endpoint exclusions}
\label{slice:early-exclusions}

We now apply the section identities to the two vertex patterns in
Proposition~\ref{geo:patterns}.  The four-vertex pattern is excluded at
both transition volumes by the same estimate.  The three-vertex pattern
is excluded at the first transition; at the second, the identities give
the parameters needed for the nesting argument.

\subsection{The four-vertex pattern}

\begin{proposition}\label{slice:four-vertex}
A nonstandard minimizer at either $V_1=4\pi/81$ or $V_2=1/\pi$
cannot have the occupied vertex pattern
$\{000,100,010,001\}$.
\end{proposition}

\begin{proof}
The complement of this pattern becomes the same pattern after reversing
all three coordinates.  We may therefore choose the smaller-volume side,
so that $V\in\{V_1,V_2\}$ and $h\ge0$.  Put
\[
 a=\sqrt{\pi V},\qquad x=hA.
\]
The standard competitors and the curvature estimate
\eqref{pre:curvature} give
\begin{equation}\label{slice:four-parameters}
 0<A\le a\le1,\qquad
 0\le x\le\sqrt{\pi A}\le\sqrt\pi.
\end{equation}
Since $b_x+b_y+b_z=2$, choose a slicing direction, denoted by $z$,
for which $b=b_z\ge2/3$, and write $d=\sqrt{1-b}$.
The strict positivity of $N_z$ in the interior implies $b<1$, so
\[
 0<d\le\frac1{\sqrt3}.
\]
The top section is a corner region, separated from the two opposite
edges of its square.  The bottom section has a complement of the same
kind at the opposite corner.  Reflecting either corner region in its
two adjacent edges gives a planar region of four times its area and
four times its relative perimeter.  The planar isoperimetric inequality
therefore gives
\[
 p(1)\ge\sqrt{\pi r},\qquad
 p(0)\ge\sqrt{\pi(1-s)}.
\]
Monotonicity gives $r\le V$.  Define $w\in[0,1]$ by
$\sqrt{\pi r}=aw$.  The endpoint flux identities
\eqref{slice:flux} and the projection estimate
\eqref{slice:projection} now imply
\begin{equation}\label{slice:four-endpoint-inequalities}
\begin{aligned}
 aw&\le Ab-\frac{ha^2}{\pi}(1-w^2),\\
 \sqrt{\pi(1-s)}
 &\le Ab-\frac{ha^2}{\pi}(1-w^2)+xd,\\
 s&\le\frac{a^2w^2}{\pi}+Ad.
\end{aligned}
\end{equation}
Indeed, $p(0)=p(1)+h(s-r)$, and $s-r\le Ad$.

Dividing the first inequality in
\eqref{slice:four-endpoint-inequalities} by $a$ and using $A\le a$
and $ha\ge x$ yields
\[
 f(w):=w+\frac{x}{\pi}(1-w^2)\le b.
\]
Since $f(1)=1>b$, continuity gives $\tau\in[w,1)$ with
\begin{equation}\label{slice:four-root}
 \tau+\frac{x}{\pi}(1-\tau^2)=1-d^2,
 \qquad
 x=\frac{\pi(1-\tau-d^2)}{1-\tau^2}.
\end{equation}
In particular $\tau\le1-d^2$.  No monotonicity of $f$ is required.
The function $E(q)=Ab-ha^2(1-q^2)/\pi$ is nondecreasing for $q\ge0$,
and
\[
 E(\tau)
 =A\tau+\frac{h(A^2-a^2)}{\pi}(1-\tau^2)
 \le a\tau.
\]
Consequently \eqref{slice:four-endpoint-inequalities} gives
\[
 \sqrt{\pi(1-s)}\le a\tau+xd,
 \qquad
 s\le\frac{a^2\tau^2}{\pi}+Ad.
\]
Both sides of the first inequality are nonnegative.  Squaring it,
combining the two inequalities, and using $A,a\le1$, we obtain the
necessary condition
\begin{equation}\label{slice:four-contradiction}
 1\le R(\tau,d):=
 d+\frac{\tau^2+(\tau+xd)^2}{\pi}.
\end{equation}
The scalar inequality in Lemma~\ref{num:four-vertex-scalar}
contradicts \eqref{slice:four-contradiction}; its hypotheses follow from
\eqref{slice:four-parameters} and \eqref{slice:four-root}, together with
$0<d\le1/\sqrt3$.  This excludes the four-vertex pattern at both
transition volumes.
\end{proof}

\subsection{The three-vertex pattern}

Consider the pattern $\{000,100,010\}$, with $z$ the missing axial
direction.  At this stage the occupied side may a priori be either
side of the original minimizer.  The top square is vacant, and is
separated from the region by a nonempty collar, so $v(z)=p(z)=0$ near
$z=1$.  The flux constant in \eqref{slice:flux} is therefore zero:
$b_zA=hV$.  The bottom section has a corner complement and
$0<s<1$.  Its planar perimeter bound and
\eqref{slice:projection} give
\begin{equation}\label{slice:three-flux}
 \int\alpha\,d\ell=hv(z),\qquad
 p(0)=hs\ge\sqrt{\pi(1-s)},\qquad
 s^2\le A(A-hV).
\end{equation}
The first identity holds at every regular slicing level.  In particular
$h>0$, since the bottom complement has positive area.  Profile concavity
and complement symmetry give nonpositive outward curvature on the
larger-volume side.  Thus this pattern necessarily labels the
smaller-volume side, and $V$ equals the endpoint volume under
consideration.

\begin{proposition}\label{slice:three-first}
A nonstandard minimizer at $V_1=4\pi/81$ cannot have the three-vertex
pattern.
\end{proposition}

\begin{proof}
Set $a=\sqrt{\pi V_1}=2\pi/9$ and $x=hA$.  As before,
$A\le a$ and $0\le x\le M:=\sqrt{\pi a}=\pi\sqrt2/3<\pi$.
Since $V_1=a^2/\pi\ge A^2/\pi$,
\eqref{slice:three-flux} gives
\[
 s^2\le A^2-xV_1\le A^2(1-x/\pi),
 \qquad s\le A\sqrt{1-x/\pi}.
\]
Combining this with $\pi(1-s)\le h^2s^2$ yields
\begin{equation}\label{slice:three-first-scalar}
 1\le a\sqrt{1-x/\pi}
       +\frac{x^2}{\pi}(1-x/\pi).
\end{equation}
The first summand decreases in $x$.  The second, denoted by $g(x)$,
has derivative $g'(x)=x(2\pi-3x)/\pi^2\ge0$ on $[0,M]$, since
$M<2\pi/3$.
For $0\le x\le1$, the right side of
\eqref{slice:three-first-scalar} is at most
\[
 a+g(1)
 <\frac{44}{63}+\frac29=\frac{58}{63}<1.
\]
Here $(\pi-1)/\pi^2<2/9$ because $(p-1)/p^2$ decreases for $p>2$
and $\pi>3$.  For $1\le x\le M$, it is at most
\begin{align*}
 a\sqrt{1-1/\pi}+g(M)
 &=a\left(\sqrt{1-1/\pi}+1-\frac{\sqrt2}{3}\right)\\
 &<\frac{44}{63}\left(\frac56+\frac8{15}\right)
 =\frac{902}{945}<1.
\end{align*}
The strict bounds $\sqrt{1-1/\pi}<5/6$ and $\sqrt2>7/5$ follow
by squaring, using $\pi<22/7$.  Both intervals contradict
\eqref{slice:three-first-scalar}.
\end{proof}

\begin{proposition}\label{slice:three-second-parameters}
Suppose that a nonstandard minimizer at $V_2=1/\pi$ has the
three-vertex pattern.  Then $A\le1$, the identities
\eqref{slice:three-flux} hold, and
\begin{equation}\label{slice:three-parameters}
 \frac{\sqrt{\pi(1-s)}}s\le h\le\pi(1-s^2),\qquad
 0.535<s<0.89,\qquad 0.64\le h\le2.25,\qquad s\ge t(h),
\end{equation}
where
\begin{equation}\label{slice:t-definition}
 t=t(h):=\frac{2}{1+\sqrt{1+4h^2/\pi}},
 \qquad h^2t^2=\pi(1-t).
\end{equation}
\end{proposition}

\begin{proof}
The slab competitor gives $A\le1$.  The lower bound on $h$ in
\eqref{slice:three-parameters} follows from the bottom perimeter
inequality in \eqref{slice:three-flux}.  Its projection inequality
gives
\[
 h\le\pi\left(A-\frac{s^2}{A}\right)
 \le\pi(1-s^2),
\]
since $A-s^2/A$ is increasing for $A>0$.  Solving
$h^2s^2+\pi s-\pi\ge0$ for $s>0$ gives $s\ge t(h)$.

Combining the two bounds on $h$, and dividing by $1-s>0$, gives
\[
 \Phi(s):=s^2(1-s)(1+s)^2\ge\frac1\pi.
\]
Its derivative is
$\Phi'(s)=s(1+s)(2+s-5s^2)$, so $\Phi$ increases up to
$s_*=(1+\sqrt{41})/10$ and decreases thereafter.
The points $107/200$ and $89/100$ lie on the respective sides of
$s_*$.  Exact arithmetic gives
\[
 \frac{22}{7}\Phi(107/200)
 =\frac{1103875107423}{1120000000000}<1,
 \qquad
 \frac{22}{7}\Phi(89/100)
 =\frac{4890924423}{5000000000}<1.
\]
Since $\pi<22/7$, unimodality implies
$107/200<s<89/100$.
Finally, $\sqrt{\pi(1-s)}/s$ decreases with $s$, and hence
\[
 h>\sqrt{\frac{3300}{7921}}>\frac{16}{25},
 \qquad
 h<\frac{22}{7}\left(1-\left(\frac{107}{200}\right)^2\right)
 =\frac{314061}{140000}<\frac94.
\]
For the first strict comparison, $\pi>3$ was used; its last step
follows from
$3300/7921-(16/25)^2=34724/4950625>0$.
These strict bounds imply the closed curvature interval stated in
\eqref{slice:three-parameters}.
\end{proof}
